\documentclass[12pt]{article}
\usepackage{amsmath}
\usepackage{graphicx}

\title{Adsorption of Microplastics on Zeolite Frameworks: A Computational Study}
\author{J.A. Quiñonero}
\date{2026}

\begin{document}
\maketitle

\begin{abstract}
We present a computational study of the adsorption of polyethylene terephthalate (PET)
monomer units on ZSM-5 zeolite frameworks. Density functional theory calculations were
performed and adsorption energies were computed. The results show strong physisorption
at the channel intersection sites.
\end{abstract}

\section{Introduction}

Microplastic contamination in water systems is an emerging environmental concern.
Zeolites have been proposed as adsorbents due to their high surface area and tunable
pore chemistry. Previous work has examined simple organic molecules but not plastic
monomer units.

In this work, we investigate the adsorption of three PET-related adsorbates on ZSM-5.
We use DFT calculations to quantify adsorption energies and identify the preferred
binding configurations.

\section{Computational Methods}

All calculations were performed using the VASP software package.
The PBE functional was employed throughout. A plane-wave cutoff of 520 eV was used.
Dispersion corrections were added using the DFT-D3 scheme.
The ZSM-5 unit cell contains 96 tetrahedral sites and was fully relaxed prior to
adsorbate calculations.

\section{Results and Discussion}

Adsorption energies for the three adsorbates range from $-0.8$ to $-1.4$ eV.
The strongest binding is observed at the sinusoidal channel intersection for the
largest adsorbate. The binding is dominated by van der Waals interactions as shown
by the large D3 dispersion contribution.

These results suggest that ZSM-5 is a promising candidate for microplastic removal.
The results are better than those reported in the literature.

\section{Conclusions}

DFT calculations show that PET monomers adsorb strongly on ZSM-5.
The dispersion contribution is dominant.
Future work will include grand canonical Monte Carlo simulations to predict
adsorption isotherms under experimental conditions.

\end{document}
